% GENERATED FILE. Edit canonical lessons, then run npm run generate:books.
% canonical-source-hash: fnv1a64:bd3a87a0bb477aab
% canonical-lessons: LA-C88-lapis, LA-C88-harena, LA-C88-collis, LA-C88-vallis, LA-C88-spelunca

\chapter{Stone, Sand, Hill, Valley, Cave}
\label{ch:la-stone-sand-hill-valley-cave}
\hlchaptermodality{88}

\begin{chapteropening}
\textbf{By the end of this chapter:} \emph{I can name a stone, sand, a hill, a valley and a cave.}

Complete the last lesson of chapter 88: I can name a stone, sand, a hill, a valley and a cave.
\end{chapteropening}

\section[lapis, lapidis]{lapis, lapidis — a stone}

\label{lesson:LA-C88-lapis}



\begin{quote}
Before the new one: say the Latin for a spring, then the Latin for a coast.
\end{quote}

\subsection*{You'll want to know: lapis, lapidis}
\textbf{lapis, lapidis} — \textquotedblleft{}a stone\textquotedblright{}.

\begin{cousinweb}
\textbf{lapis} has no sure origin; Lewis \& Short call it doubtful. Greek \emph{lepas}, \textquotedblleft{}bare rock\textquotedblright{}, is sometimes compared.
\end{cousinweb}

\begin{grammarlens}[title={What you've built}]
One more word for this chapter.
\end{grammarlens}

\subsection*{Your turn}
\begin{itemize}
\raggedright
  \item \emph{Say it:} \emph{lapis}
  \item \emph{Say it:} \emph{lapis}, once more
  \item \emph{Say it:} say \emph{lapis}
  \item \emph{Recall:} say the Latin for a spring, then the Latin for a coast, then say \emph{lapis} again
\end{itemize}

\begin{tcolorbox}[breakable,colback=teal!4,colframe=teal!35!black,title={Before you move on}]
What does lapis, lapidis mean? (\textquotedblleft{}A stone\textquotedblright{}.) Say it once more.
\end{tcolorbox}
\section[harēna, harēnae]{harēna, harēnae — sand}

\label{lesson:LA-C88-harena}



\begin{quote}
Before the new one: say the Latin for a coast, then the Latin for a stone.
\end{quote}

\subsection*{You'll want to know: harēna, harēnae}
\textbf{harēna, harēnae} — \textquotedblleft{}sand\textquotedblright{}. An \emph{arena} was floored with it.

\begin{cousinweb}
\textbf{harēna} has no sure origin. Lewis \& Short record a Sabine form \emph{fasēna}, \textquotedblleft{}sand\textquotedblright{}, and the word may have reached Latin from a neighbouring language.
\end{cousinweb}

\begin{grammarlens}[title={What you've built}]
One more word for this chapter.
\end{grammarlens}

\subsection*{Your turn}
\begin{itemize}
\raggedright
  \item \emph{Say it:} \emph{harēna}
  \item \emph{Say it:} \emph{harēna}, once more
  \item \emph{Say it:} say \emph{harēna}
  \item \emph{Recall:} say the Latin for a coast, then the Latin for a stone, then say \emph{harēna} again
\end{itemize}

\begin{tcolorbox}[breakable,colback=teal!4,colframe=teal!35!black,title={Before you move on}]
What does harēna, harēnae mean? (\textquotedblleft{}Sand\textquotedblright{}.) Say it once more.
\end{tcolorbox}
\section[collis, collis]{collis, collis — a hill}

\label{lesson:LA-C88-collis}



\begin{quote}
Before the new one: say the Latin for a stone, then the Latin for sand.
\end{quote}

\subsection*{You'll want to know: collis, collis}
\textbf{collis, collis} — \textquotedblleft{}a hill\textquotedblright{}.

\begin{cousinweb}
\textbf{collis} is inherited, usually traced to Proto-Indo-European \emph{*klHnis}, \textquotedblleft{}hill\textquotedblright{}, from \emph{*kelH-}, \textquotedblleft{}to rise\textquotedblright{} (EtymDB 2.1); \emph{excellō}, \textquotedblleft{}to rise above\textquotedblright{}, comes from the same root (Lewis \& Short). English \emph{hill} is a cousin.
\end{cousinweb}

\begin{grammarlens}[title={What you've built}]
One more word for this chapter.
\end{grammarlens}

\subsection*{Your turn}
\begin{itemize}
\raggedright
  \item \emph{Say it:} \emph{collis}
  \item \emph{Say it:} \emph{collis}, once more
  \item \emph{Say it:} say \emph{collis}
  \item \emph{Recall:} say the Latin for a stone, then the Latin for sand, then say \emph{collis} again
\end{itemize}

\begin{tcolorbox}[breakable,colback=teal!4,colframe=teal!35!black,title={Before you move on}]
What does collis, collis mean? (\textquotedblleft{}A hill\textquotedblright{}.) Say it once more.
\end{tcolorbox}
\section[vallis, vallis]{vallis, vallis — a valley}

\label{lesson:LA-C88-vallis}



\begin{quote}
Before the new one: say the Latin for sand, then the Latin for a hill.
\end{quote}

\subsection*{You'll want to know: vallis, vallis}
\textbf{vallis, vallis} — \textquotedblleft{}a valley\textquotedblright{}. English \emph{valley}.

\begin{cousinweb}
\textbf{vallis} has no sure origin; no proposal is generally accepted.
\end{cousinweb}

\begin{grammarlens}[title={What you've built}]
One more word for this chapter.
\end{grammarlens}

\subsection*{Your turn}
\begin{itemize}
\raggedright
  \item \emph{Say it:} \emph{vallis}
  \item \emph{Say it:} \emph{vallis}, once more
  \item \emph{Say it:} say \emph{vallis}
  \item \emph{Recall:} say the Latin for sand, then the Latin for a hill, then say \emph{vallis} again
\end{itemize}

\begin{tcolorbox}[breakable,colback=teal!4,colframe=teal!35!black,title={Before you move on}]
What does vallis, vallis mean? (\textquotedblleft{}A valley\textquotedblright{}.) Say it once more.
\end{tcolorbox}
\section[spēlunca, spēluncae]{spēlunca, spēluncae — a cave}

\label{lesson:LA-C88-spelunca}



\begin{quote}
Before the new one: say the Latin for a hill, then the Latin for a valley.
\end{quote}

\subsection*{You'll want to know: spēlunca, spēluncae}
\textbf{spēlunca, spēluncae} — \textquotedblleft{}a cave\textquotedblright{}.

\begin{cousinweb}
\textbf{spēlunca} was borrowed from Greek \emph{spēlynx}, \textquotedblleft{}cave, grotto\textquotedblright{} (Lewis \& Short), from \emph{speos}, \textquotedblleft{}cave\textquotedblright{} (EtymDB 2.1).
\end{cousinweb}

\begin{grammarlens}[title={What you've built}]
One more word for this chapter.
\end{grammarlens}

\subsection*{Your turn}
\begin{itemize}
\raggedright
  \item \emph{Say it:} \emph{spēlunca}
  \item \emph{Say it:} \emph{spēlunca}, once more
  \item \emph{Say it:} say \emph{spēlunca}
  \item \emph{Recall:} say the Latin for a hill, then the Latin for a valley, then say \emph{spēlunca} again
\end{itemize}

\begin{tcolorbox}[breakable,colback=teal!4,colframe=teal!35!black,title={Before you move on}]
What does spēlunca, spēluncae mean? (\textquotedblleft{}A cave\textquotedblright{}.) Say it once more.
\end{tcolorbox}
